% TOP_PROBLEM: {"id": 1100118, "title": "Questions in Geometric Group Theory — Q 1.18", "category_id": 17, "category": {"id": 17, "name": "group_theory", "display_name": "Group Theory", "description": "Problems about groups, group actions, representations, and related algebraic structures.", "slug": "group-theory", "order_index": 17, "created_at": "2026-07-31T15:26:25.671Z"}, "status": "partially_solved", "problem_number": "AMR-010-0118", "set_id": 13, "statusQualification": "The group-ring comparison uses reduced boundary cohomology. The sphere-recognition subquestion has a precise yes/no output."}
% FIRST_POSTED: 2026-10-02
% ENTRY_KIND: statement_only
% UnsolvedMath ID 1100118; source catalogue code AMR-010-0118
% !TeX program = lualatex
% Definition and sources only; this file is not an LLM solution attempt.
\documentclass[11pt,a4paper]{article}
\usepackage[margin=25mm]{geometry}
\usepackage{fontspec}
\setmainfont{lmroman10-regular.otf}[BoldFont=lmroman10-bold.otf,
 ItalicFont=lmroman10-italic.otf,BoldItalicFont=lmroman10-bolditalic.otf]
\usepackage{amsmath,amssymb,mathtools}
\usepackage{unicode-math}
\setmathfont{latinmodern-math.otf}
\AtBeginDocument{\let\setminus\smallsetminus}
\usepackage{xurl,hyperref}
\hypersetup{colorlinks=true,urlcolor=blue}
\setcounter{secnumdepth}{0}
\setlength{\parindent}{0pt}
\setlength{\parskip}{5pt}
\setlength{\emergencystretch}{3em}
\begin{document}
\section*{Questions in Geometric Group Theory — Q 1.18}\label{1100118}
{\small UnsolvedMath ID 1100118 · AMR-010-0118 · Group Theory · AMR Open Problem Lists}
\subsection{Definitions and mathematical statement}
Let $G$ be given by a finite presentation, with the promise that it is word-hyperbolic: its Cayley graph has uniformly thin geodesic triangles. Its Gromov boundary $\partial G$ is the compact space of geodesic rays from a basepoint, with rays identified when they remain a bounded distance apart. Write $\check H^i(\partial G;\mathbb Z)$ for integral \v{C}ech cohomology, computed as the direct limit of the cohomology of nerves of finite open covers under refinement.

Find an algorithm that determines these cohomology groups from the presentation. A concrete decision problem, avoiding an unspecified output format for arbitrary abelian groups, is whether there is an algorithm which always halts and decides
\[
 \check H^i(\partial G;\mathbb Z)\cong
 \check H^i(S^2;\mathbb Z)\quad\hbox{for every }i\geq0.
\]
Thus the unreduced groups sought are $\mathbb Z$ in degrees zero and two, and zero in all other degrees. In reduced notation the only nonzero group is $\widetilde{\check H}^{2}\cong\mathbb Z$.

The comparison with group-ring cohomology in the source uses reduced boundary cohomology; retaining that convention avoids a spurious degree-zero term. The input is finite algebraic data, not an already supplied triangulation of the boundary. The boundary can be a nonmanifold compactum, so computing ordinary simplicial cohomology of a single finite complex is not assumed to suffice.
\subsection{Short English statement}
Can one algorithmically compute boundary cohomology, and in particular recognize the integral cohomology of the two-sphere?
\subsection{Sources}
\begin{itemize}
\item[S1] ULAM.AI and the UnsolvedMath Contributors, supplied problems.json, record 1100118 (AMR-010-0118), AMR Open Problem Lists. Catalogue identity and source transcription; CC BY 4.0.
\url{https://huggingface.co/datasets/ulamai/UnsolvedMath}.
\item[S2] Bestvina - Questions in Geometric Group Theory (pdf) (2004), Question 1.18 (PDF page 4). Original source indexed by AMR-010-0118.
\url{https://www.math.utah.edu/~bestvina/eprints/questions-updated.pdf}.
\end{itemize}
\end{document}
